\section{Exploratory Crossover Study}
\label{sec:study}

\subsection{Protocol}

The implemented study asks whether the relative adaptive policy feels more
appropriate than Random selection from the same catalog. Every participant
completes one five-track Random block and one five-track Vibe block. Two masked
protocols reverse the order, and odd/even participant numbers suggest the
assignment. Participants see a protocol identifier rather than the condition
name when the operator keeps the participant-facing mode enabled. The same
interface exposes a view toggle and protocol controls, while the experimenter
and exported data can reveal the order. Masking therefore depended on study
administration; it was not technically enforced or double blind. The paired
design is efficient but remains susceptible to practice and carryover, making
reversed order a substantive control~\cite{greenwald1976}.

Each track plays for up to 60 seconds. Participants may select \emph{Rate now};
the actual duration and an early-rating flag are preserved. After each track,
the interface collects a seven-point liking rating, a seven-point mood-fit
rating, and a categorical current-mood report. Mood fit is treated as the
primary outcome in this exploratory report because it directly evaluates the
matching idea. Liking is treated as secondary and helps separate contextual fit
from general song preference.

\subsection{Available evidence}

The supplied analysis report publishes one five-trial condition mean per
participant/session for 15 participants, together with paired tests and
condition-level summaries. These 15 rows are frozen in the repository under
generic analysis IDs. The protocol implies 75 planned ratings per condition and
outcome, but the repository still lacks the underlying trial rows needed to
verify completeness, exclusions, or within-session variation. It also lacks
recruitment records, demographics, and the complete study-administration file.
The present paper therefore treats the participant/session as the inferential
unit and does not claim a fully reproducible human-study archive.

Paired differences are defined as Vibe minus Random. We show the published
participant/session means, report condition means and standard deviations,
paired mean differences with two-sided 95\% confidence intervals, and the
originally supplied paired tests. Because no dated preregistration was
available, all inference is exploratory~\cite{nosek2018}. We foreground paired
effects and confidence intervals rather than a thresholded
$p$-value~\cite{cumming2014}.

\begin{table}[t]
  \caption{Exploratory condition summaries ($N=15$ participants). Confidence
  intervals are two-sided. The reported study tests used one-sided
  Vibe-over-Random alternatives.}
  \label{tab:pilot}
  \centering
  \small
  \begin{tabular}{@{}lccc@{}}
    \toprule
    Outcome & Vibe $M$ (SD) & Random $M$ (SD) & Difference [95\% CI] \\
    \midrule
    Mood fit & 4.71 (0.99) & 4.15 (1.30) & $+0.56$ [$-0.21,1.33$] \\
    Liking   & 4.56 (1.18) & 4.47 (0.77) & $+0.09$ [$-0.52,0.71$] \\
    \bottomrule
  \end{tabular}
\end{table}
